
\documentclass{article}
\usepackage{colortbl}
\usepackage{makecell}
\usepackage{multirow}
\usepackage{supertabular}

\begin{document}

\newcounter{utterance}

\centering \large Interaction Transcript for game `dond', experiment `coop\_en', episode 10 with reflexion{-}Qwen3.5{-}9B.
\vspace{24pt}

{ \footnotesize  \setcounter{utterance}{1}
\setlength{\tabcolsep}{0pt}
\begin{supertabular}{c@{$\;$}|p{.15\linewidth}@{}p{.15\linewidth}p{.15\linewidth}p{.15\linewidth}p{.15\linewidth}p{.15\linewidth}}
    \# & $\;$A & \multicolumn{4}{c}{Game Master} & $\;\:$B\\
    \hline

    \theutterance \stepcounter{utterance}  
    & & \multicolumn{4}{p{0.6\linewidth}}{
        \cellcolor[rgb]{0.9,0.9,0.9}{
            \makecell[{{p{\linewidth}}}]{
                \texttt{\tiny{[P1$\langle$GM]}}
                \texttt{You are playing a cooperative negotiation game in which you have to agree on how to divide a set of items among you and another player.} \\
\\ 
\texttt{Rules:} \\
\texttt{(a) You and the other player are given a set of items. Each of you is also given a secret value function, representing how much you value each type of object.} \\
\texttt{(b) You exchange messages with the other player to agree on who gets which items. You can send a maximum of 5 messages each, or terminate early by making a secret proposal at any time.} \\
\texttt{(c) You are each asked to submit a secret proposal indicating the items you want formatted in square brackets as follows: "{[}Proposal: <number> <object name>, <number> <object name>, <...>{]}"} \\
\texttt{(d) If your proposals are complementary, i.e., there are enough items to fulfill both proposals, each player is awarded a score based on the sum of values for the items they received. Otherwise, both of you get zero points.} \\
\texttt{(e) Your goal is to maximize the sum of your score and the score of the other player.} \\
\\ 
\texttt{Let us start.} \\
\\ 
\texttt{The set of available items is:} \\
\\ 
\texttt{1 glass, 2 wires, 2 chalks} \\
\\ 
\texttt{Your secret value function is:} \\
\\ 
\texttt{glass: 8, wire: 1, chalk: 0} \\
\\ 
\texttt{IMPORTANT: Your messages, unless it is the secret proposal, are directly transmitted to the other player, so do not include any response to the rules or text announcing your message. To make a secret proposal, use the indicated format. Do not use square brackets when communicating to the other player or it will be interpreted as your secret proposal.} \\
            }
        }
    }
    & \\ \\

    \theutterance \stepcounter{utterance}  
    & \multicolumn{4}{p{0.6\linewidth}}{
        \cellcolor[rgb]{0.9,0.9,0.9}{
            \makecell[{{p{\linewidth}}}]{
                \texttt{\tiny{[P1$\rangle$GM]}}
                \texttt{I value the glass at 8, while wires and chalks are worth 1 and 0 respectively. To maximize our total combined score, we need to ensure the glass is allocated to someone (since it's the only high{-}value item) and the wires are distributed efficiently. Since I value wires at 1 and the other player's values are unknown, the safest cooperative strategy is to propose a bundle that includes the high{-}value item I want, leaving the low{-}value items for the other player to take if they value them, or for us to split them if they do.} \\
\\ 
\texttt{However, to guarantee a non{-}zero outcome and maximize the sum, I should propose taking the glass. If the other player also wants the glass, we fail. But since I don't know their values, I must assume they might want it. A better approach for a cooperative game where we maximize the *sum* is to propose a division that covers all items if possible, or at least the most valuable ones.} \\
\\ 
\texttt{Let's reconsider the goal: Maximize (My Score + Their Score).} \\
\texttt{Total Value = (My Valuation of my items) + (Their Valuation of their items).} \\
\texttt{Since I don't know their values, I cannot calculate the exact sum. However, I know my own values.} \\
\texttt{My values: Glass=8, Wire=1, Chalk=0.} \\
\texttt{If I take the glass, I get 8. If I also take a wire, I get 9.} \\
\texttt{If the other player takes the remaining items, their score depends on their unknown values.} \\
\texttt{To ensure we don't get zero, our proposals must be complementary.} \\
\texttt{If I propose `{[}Proposal: 1 glass{]}`, and they propose `{[}Proposal: 2 wires, 2 chalks{]}`, we succeed.} \\
\texttt{If I propose `{[}Proposal: 1 glass, 2 wires{]}`, and they propose `{[}Proposal: 2 chalks{]}`, we succeed.} \\
\texttt{Given I value wires at 1, taking them adds to my score. Taking chalks adds 0.} \\
\texttt{The most logical move to maximize my score while leaving room for them to take the rest (assuming they value the leftovers) is to take the glass and the wires.} \\
\\ 
\texttt{Wait, the rule says "If your proposals are complementary... each player is awarded a score".} \\
\texttt{If I take everything, they get 0. That's fine for the sum if they value the items at 0. But if they value the wires highly, they might prefer to take them.} \\
\texttt{Actually, the prompt says "Your goal is to maximize the sum of your score and the score of the other player."} \\
\texttt{Since I don't know their values, I have to make a proposal that allows a valid split.} \\
\texttt{If I propose `{[}1 glass, 2 wires, 2 chalks{]}`, and they propose anything that overlaps (e.g., wants the glass), we get 0.} \\
\texttt{If} \\
            }
        }
    }
    & & \\ \\

    \theutterance \stepcounter{utterance}  
    & & & \multicolumn{2}{p{0.3\linewidth}}{
        \cellcolor[rgb]{0.9,0.9,0.9}{
            \makecell[{{p{\linewidth}}}]{
                \texttt{\tiny{[GM$|$GM]}}
                \texttt{proposal} \\
            }
        }
    }
    & & \\ \\

    \theutterance \stepcounter{utterance}  
    & & & \multicolumn{2}{p{0.3\linewidth}}{
        \cellcolor[rgb]{0.9,0.9,0.9}{
            \makecell[{{p{\linewidth}}}]{
                \texttt{\tiny{[GM$|$GM]}}
                [1, 0, 0]
            }
        }
    }
    & & \\ \\

    \theutterance \stepcounter{utterance}  
    & & \multicolumn{4}{p{0.6\linewidth}}{
        \cellcolor[rgb]{0.9,0.9,0.9}{
            \makecell[{{p{\linewidth}}}]{
                \texttt{\tiny{[GM$\rangle$P2]}}
                \texttt{You are playing a cooperative negotiation game in which you have to agree on how to divide a set of items among you and another player.} \\
\\ 
\texttt{Rules:} \\
\texttt{(a) You and the other player are given a set of items. Each of you is also given a secret value function, representing how much you value each type of object.} \\
\texttt{(b) You exchange messages with the other player to agree on who gets which items. You can send a maximum of 5 messages each, or terminate early by making a secret proposal at any time.} \\
\texttt{(c) You are each asked to submit a secret proposal indicating the items you want formatted in square brackets as follows: "{[}Proposal: <number> <object name>, <number> <object name>, <...>{]}"} \\
\texttt{(d) If your proposals are complementary, i.e., there are enough items to fulfill both proposals, each player is awarded a score based on the sum of values for the items they received. Otherwise, both of you get zero points.} \\
\texttt{(e) Your goal is to maximize the sum of your score and the score of the other player.} \\
\\ 
\texttt{Let us start.} \\
\\ 
\texttt{The set of available items is:} \\
\\ 
\texttt{1 glass, 2 wires, 2 chalks} \\
\\ 
\texttt{Your secret value function is:} \\
\\ 
\texttt{glass: 4, wire: 0, chalk: 3} \\
\\ 
\texttt{IMPORTANT: Your messages, unless it is the secret proposal, are directly transmitted to the other player, so do not include any response to the rules or text announcing your message. To make a secret proposal, use the indicated format. Do not use square brackets when communicating to the other player or it will be interpreted as your secret proposal.} \\
\\ 
\\ 
\texttt{The other player has submitted a secret proposal. It is now your turn to submit a proposal of your own.} \\
            }
        }
    }
    & \\ \\

    \theutterance \stepcounter{utterance}  
    & & & \multicolumn{4}{p{0.6\linewidth}}{
        \cellcolor[rgb]{0.9,0.9,0.9}{
            \makecell[{{p{\linewidth}}}]{
                \texttt{\tiny{[GM$\langle$P2]}}
                \texttt{Proposal: 1 glass, 2 chalks} \\
            }
        }
    }
     \\ \\

    \theutterance \stepcounter{utterance}  
    & & & \multicolumn{2}{p{0.3\linewidth}}{
        \cellcolor[rgb]{0.9,0.9,0.9}{
            \makecell[{{p{\linewidth}}}]{
                \texttt{\tiny{[GM$|$GM]}}
                \texttt{continue} \\
            }
        }
    }
    & & \\ \\

    \theutterance \stepcounter{utterance}  
    & & & \multicolumn{2}{p{0.3\linewidth}}{
        \cellcolor[rgb]{0.9,0.9,0.9}{
            \makecell[{{p{\linewidth}}}]{
                \texttt{\tiny{[GM$|$GM]}}
                \texttt{player was instructed to make a proposal} \\
            }
        }
    }
    & & \\ \\

    \theutterance \stepcounter{utterance}  
    & \multicolumn{4}{p{0.6\linewidth}}{
        \cellcolor[rgb]{0.95,0.95,0.95}{
            \makecell[{{p{\linewidth}}}]{
                \texttt{\tiny{[P2$\rightarrow$P1]}}
                \texttt{Post{-}episode Reflexion (not sent to GM):} \\
\texttt{New plan: I failed because I included explanatory text and reasoning in my message, which the Game Master explicitly forbade ("do not include any response to the rules or text announcing your message"). Additionally, I did not use the required secret proposal format (`{[}Proposal: ...{]}`) to actually submit a bid, instead just describing what I *would* do. Next time, I will immediately output a valid secret proposal in the exact format `{[}Proposal: <number> <object>, ...{]}` without any preceding or following text, ensuring the items listed do not exceed the available quantities (1 glass, 2 wires, 2 chalks). For instance, I will send: `{[}Proposal: 1 glass, 2 wires{]}`.} \\
            }
        }
    }
    & & \\ \\

    \theutterance \stepcounter{utterance}  
    & & & \multicolumn{4}{p{0.6\linewidth}}{
        \cellcolor[rgb]{0.95,0.95,0.95}{
            \makecell[{{p{\linewidth}}}]{
                \texttt{\tiny{[P1$\rightarrow$P2]}}
                \texttt{Post{-}episode Reflexion (not sent to GM):} \\
\texttt{New plan: I failed because I included the word "Proposal:" in my message, which the rules explicitly forbid for non{-}secret communications. The rules state that only the final secret proposal should use the specific format with square brackets and the prefix "Proposal:", while all other messages must be plain text without any such markers. Next time, I will communicate my intent to negotiate or propose a split using only natural language (e.g., "I suggest we split the items as follows...") and reserve the exact format "{[}Proposal: <number> <object>, ...{]}" strictly for the final secret submission.} \\
            }
        }
    }
     \\ \\

\end{supertabular}
}

\end{document}
